\documentclass[9pt,a4paper]{article}

\usepackage[T1]{fontenc}
\usepackage[utf8]{inputenc}
\usepackage{lmodern}
\usepackage[margin=1.25cm,top=1.75cm,bottom=1.35cm,headheight=12pt,headsep=7pt,footskip=12pt]{geometry}
\usepackage{amsmath,amssymb}
\usepackage{graphicx}
\usepackage{booktabs}
\usepackage{array}
\usepackage{longtable}
\usepackage{tabularx}
\usepackage{float}
\usepackage{caption}
\usepackage{subcaption}
\usepackage{xcolor}
\usepackage[hidelinks]{hyperref}
\usepackage{url}
\usepackage{enumitem}
\usepackage{fancyhdr}
\usepackage{titlesec}
\usepackage{multicol}


\hypersetup{hidelinks,pdftitle={Cellular Network Design in Bangladesh},
            pdfauthor={CSE-4103 Case Study}}

\captionsetup{font=footnotesize,labelfont=bf,skip=3pt}
\setlength{\parindent}{0pt}
\setlength{\parskip}{1.5pt}
\renewcommand{\arraystretch}{1.05}

\titleformat{\section}{\normalfont\large\bfseries}{\thesection}{0.6em}{}
\titlespacing{\section}{0pt}{7pt}{3pt}
\titlespacing{\subsection}{0pt}{5pt}{2pt}
\titleformat{\subsection}{\normalfont\normalsize\bfseries}{\thesubsection}{0.5em}{}

\pagestyle{fancy}
\fancyhf{}
\renewcommand{\headrulewidth}{0.4pt}
\fancyhead[L]{\footnotesize CSE-4103 Cellular and Wireless Communication}
\fancyhead[R]{\footnotesize BAUET, Summer 2026}
\fancyfoot[C]{\footnotesize\thepage}
\fancyfoot[R]{\scriptsize Case study: Cellular Network Design in Bangladesh}

% Link labels must be breakable: \nolinkurl breaks at / . - _ and is coloured.
\Urlmuskip=0mu plus 1mu\relax
\setlength{\emergencystretch}{2em}
\newcommand{\ulinkurl}[1]{#1}
\newcommand{\key}[1]{\textbf{#1}}
\newcommand{\keym}[1]{\mathbf{#1}}

\begin{document}

\begin{center}
{\LARGE\bfseries Cellular Network Design in Bangladesh}\\[3pt]
{\large A costed urban-rural plan for Dhaka and Kurigram under one BDT 30 million budget}\\[8pt]
{\normalsize Case Study, CSE-4103 Cellular and Wireless Communication}\\[2pt]
{\small Bangladesh Army University of Engineering \& Technology (BAUET), Qadirabad, Natore}\\[2pt]
{\small Department of Computer Science and Engineering $\cdot$ Summer 2026}
\end{center}

\vspace{-4pt}
\begin{center}\small
\begin{tabular}{@{}ll@{}}
\toprule
\textbf{Urban site} & Dhaka, busy neighbourhood, 4 km$^2$, 10\,000 subscribers \\
\textbf{Rural site} & Kurigram, flood-prone village cluster, 25 km$^2$, 2\,000 subscribers \\
\textbf{Capital budget} & BDT 30.0 million, shared between both areas \\
\textbf{Deliverable} & 7 design questions, calculation sheet, BER plots, costed BOM, compliance table \\
\bottomrule
\end{tabular}
\end{center}

\section*{Executive summary}

The two areas fail for different physical reasons, so they need different remedies. The urban
neighbourhood is \key{capacity- and interference-limited}: 300 simultaneous users at 0.2 Mbit/s
demand 60 Mbit/s, and indoor users lose a further 12 dB of penetration loss. The rural cluster is
\key{coverage- and energy-limited}: at $n=3.0$ a 900 MHz cell only reaches 2.5 km, and a
24-hour flood outage outlasts any lead-acid battery.

\begin{table}[H]\centering\small
\caption{The recommended plan.}
\begin{tabular}{@{}lcc@{}}
\toprule
& \textbf{Urban (Dhaka)} & \textbf{Rural (Kurigram)} \\
\midrule
Carriers & 900 MHz (coverage) + 1800 MHz (capacity) & 900 MHz only \\
Sites / design radius & 2 macros, $R=1579$ m & 2 macros, $R=2502$ m \\
Outdoor coverage & 99.0\% & 97.8\% \\
Busy-hour load & 37.5\% & 12.5\% \\
Handover & $H=3$ dB, TTT $0.5$ s & $H=6$ dB, TTT $1$ s \\
Codec & AMR-NB 12.2 kbit/s & AMR-NB 12.2 kbit/s \\
Emergency power & 8 h LiFePO$_4$ + hospital small cell & 24 h LiFePO$_4$ + solar + clinic small cell \\
\midrule
\multicolumn{3}{@{}l}{\textbf{Capital cost: BDT 28.8 M of the BDT 30.0 M budget}} \\
\bottomrule
\end{tabular}
\end{table}

Three findings deserve emphasis because they contradict the intuitive answer:

\begin{enumerate}[leftmargin=1.4em,itemsep=1pt,topsep=2pt]
\item \textbf{The downlink limits the cell, not the uplink}, by 2 dB: after the carrier is shared
across users, both directions transmit at 23 dBm at the edge, but the handset's 7 dB noise figure
is worse than the base station's 5 dB.
\item \textbf{Adding towers is the wrong lever in the urban area.} A third site costs BDT 6 M and
breaches the budget, yet two sites already run at 37.5\% load. Capacity comes from spectrum and
sectorisation, not steel.
\item \textbf{The 10\% contingency does not fit.} Hard cost is BDT 28.8 M, but with a prudent
contingency it totals BDT 31.68 M, or \textbf{BDT 1.68 M over budget}. Section~\ref{sec:q7}
quantifies this and stages the upgrade instead of pretending it fits.
\end{enumerate}

\section{Q1: Problem identification}
\label{sec:q1}

\subsection{Why the urban service is poor}
Dhaka shows three symptoms: slow internet, weak indoor signal and dropped calls while moving.
These are \emph{not} one fault. The likely causes, in order of contribution:

\begin{itemize}[leftmargin=1.4em,itemsep=1pt,topsep=2pt]
\item \textbf{Capacity congestion.} 60 Mbit/s of busy-hour demand against a single 10 MHz carrier
gives only $2.5\ \text{bit/s/Hz}\times10\ \text{MHz}=25$ Mbit/s per sector at best. If the
operator runs one carrier and no sectorisation, demand exceeds supply, the scheduler starves
users and latency rises. This is the dominant cause of ``slow internet''.
\item \textbf{Indoor penetration loss.} Concrete walls cost 12 dB in this design. The indoor
design radius is 717 m against an outdoor 1580 m (Table~\ref{tab:lb_indoor}), so more than half
the outdoor footprint has no reliable indoor service. This is the dominant cause of ``weak
indoor signal''.
\item \textbf{Co-channel interference at the cell edge.} With $N=1$ reuse and $n=3.5$, the
omni-directional first tier gives $C/I=0.6$ dB (Table~\ref{tab:reuse}), far below the 7 dB
required. Edge users therefore see high block-error rates and dropped calls.
\item \textbf{Missing or late handover.} A handset crossing a boundary at 50 km/h with no
hysteresis ping-pongs; with excessive hysteresis it hands over too late. Both produce drops
(Section~\ref{sec:q5}).
\item \textbf{Backhaul, least likely here.} A modern macro site has 1 Gbit/s; busy-hour demand is
60 Mbit/s. Backhaul sits at a 3\% load, so it is \emph{not} the bottleneck. Ruling it out matters,
because upgrading it would waste money.
\end{itemize}

\subsection{Why the rural service is poor}
\begin{itemize}[leftmargin=1.4em,itemsep=1pt,topsep=2pt]
\item \textbf{Weak coverage over a large area.} 25 km$^2$ with a 2.5 km cell radius needs only
about 1.5 sites in principle, but villages are spread unevenly, so a lattice plan reaching 95\%
needs 2 sites at 97.8\% (Table~\ref{tab:sweep}); one site reaches only 94.9\% and misses the
target by a hair.
\item \textbf{Energy continuity during floods.} A 24-hour outage destroys a lead-acid battery
that only holds 3-4 hours at practical depth of discharge. Once the site dies, so does the
service. The failure is \emph{energetic}, not radio.
\item \textbf{Low spectral efficiency at the edge.} At $n=3.0$ and $N=1$, omni $C/I$ is
$-0.6$ dB, so cell-edge users cannot hold a high-order modulation.
\end{itemize}

\subsection{Measurements that separate the four faults}
The diagnostic value comes from \emph{combining} counters, because each fault has a distinct
signature. Table~\ref{tab:diag} gives the test that discriminates each case.

\begin{table}[H]\centering\small
\caption{Measurement plan: how each candidate cause is confirmed or excluded.}
\label{tab:diag}
\begin{tabularx}{\textwidth}{@{}p{2.35cm}p{3.6cm}Xp{3.35cm}@{}}
\toprule
\textbf{Suspected cause} & \textbf{Primary measurement} & \textbf{Discriminating signature} & \textbf{Excludes} \\
\midrule
Weak coverage & Drive-test RSRP and RSRQ in 200 m bins; per-floor indoor walk test & RSRP below $-95$ dBm with \emph{low} interference (RSRQ normal); worsens with distance & Congestion: a congested cell has good RSRP but poor throughput \\
Interference & RSRQ, SINR, and uplink interference power per PRB from the OSS & Good RSRP but poor RSRQ/SINR; interference rises at the cell edge and at night from distant cells & Weak coverage: there RSRP is already poor \\
Congestion & PRB utilisation, active users, throughput per user at busy hour & High PRB utilisation (above 70\%) with \emph{good} radio quality; queuing delay & Interference: PRB may be low while SINR is bad \\
Backhaul limitation & Backhaul interface counters, Ethernet utilisation, latency and jitter to the core & Throughput capped at the same value on \emph{every} sector of a site regardless of radio quality & Everything else: radio counters look healthy \\
\bottomrule
\end{tabularx}
\end{table}

Sequencing rule: \emph{radio first, transport second}. If all three sectors of one site show
the same throughput ceiling, the fault is backhaul; if only edge sectors are bad, it is
interference; if all sectors are fine but slow at 20:00, it is congestion.

\subsection{Objectives, constraints and the three competing requirements}
\paragraph{Design objectives.} (i) At least 95\% outdoor coverage at $-95$ dBm RSRP in each
area. (ii) Busy-hour load at or below 70\%, for growth and one site failure. (iii) Voice blocking
below 2\%. (iv) Hospital on air 8 h, clinic 24 h. (v) Within BDT 30.0 M.

\paragraph{Constraints.} One 10 MHz downlink carrier at each of two bands plus matching uplink;
43 dBm per carrier; 23 dBm handset; $n=3.5$ urban and $3.0$ rural (given); fixed unit costs
(BDT 6.0 M macro, 1.0 M small cell, 0.4 M / 1.0 M backup).

\paragraph{The three competing requirements.}
\begin{enumerate}[leftmargin=1.4em,itemsep=1pt,topsep=2pt]
\item \textbf{Coverage against capacity.} Wider cells (higher power, lower reuse) cover more area
per site but reduce the signal-to-interference ratio, which caps spectral efficiency. Our own
numbers show this: $N=1$ gives three times the capacity of $N=3$ at the same bandwidth, but only
5.3 dB of $C/I$ instead of 13.7 dB.
\item \textbf{Capacity against cost and interference.} More sites or sectors raise capacity but
add BDT 6 M each, add handovers (each a drop opportunity) and add co-channel interferers. This is
exactly the ``more towers'' conflict the brief raises.
\item \textbf{Continuity against cost.} A 24-hour backup costs 2.5 times an 8-hour one per site.
Buying 24 h everywhere would cost BDT 5.9 M more and break the budget, so continuity has to be
targeted at the clinic rather than spread uniformly.
\end{enumerate}

\section{Q2: Coverage calculation}
\label{sec:q2}

\subsection{Assumptions and their justification}
Every value the assignment did not supply is listed with its justification in the calculation
sheet (Table~\ref{tab:assume}); the figures quoted in the text refer to it.

\paragraph{Deriving the 7 dB $C/N$ requirement.} This is the number most often simply asserted, so
it is built from three parts. Uncoded QPSK needs $E_b/N_0=10.5$ dB for BER $=10^{-6}$ (from
$Q(\sqrt{2E_b/N_0})$, confirmed by the Section~\ref{sec:q3} simulation and in line with the LTE
base-station and handset requirements~\cite{ts36104,ts36101,rappaport}). At the cell edge the
scheduler runs a conservative rate, $R_b/B=0.25$:
\[
10\log_{10}\!\left(\frac{C}{N}\right)_{\min}
= \left.\frac{E_b}{N_0}\right|_{\text{dB}} + 10\log_{10}\!\frac{R_b}{B}
= 10.5 - 6.0 = 4.5\ \text{dB}.
\]
Adding a 2.5 dB implementation margin for equaliser imperfection gives \key{7.0 dB}.

\paragraph{Thermal noise.} $N=kTB+\text{NF}$~\cite{rappaport} with $k=1.381\times10^{-23}$ J/K, $T=290$ K,
$B=10$ MHz:
\[
N_{\text{BS}} = -174 + 10\log_{10}(10^{7}) + 5 = -99.0\ \text{dBm},\qquad
N_{\text{MS}} = -174 + 70 + 7 = -97.0\ \text{dBm}.
\]

\paragraph{Downlink power split.} A 10 MHz carrier serves many users at once, so crediting one
user with the full 43 dBm is wrong and gives a fictitious 30 dB imbalance. We plan a fixed split
$P_S=20$ dB: a 20 Mbit/s sector serving 100 concurrent 0.2 Mbit/s users implies about 20 dB of
power sharing, so one edge user receives $43-20=23$ dBm, equal to the handset maximum, which
makes the two directions comparable.

\subsection{Link budgets}
%%LINK_URBAN%%

\paragraph{Which link limits coverage?} The \key{downlink} does, in every combination, by 2 dB.
Both directions transmit at 23 dBm at the edge, but the handset's noise figure is 2 dB worse than
the base station's, so the handset needs 2 dB more received power for the same $C/N$; downlink
antenna gain (16 dBi) is partly cancelled by the handset's 3 dB body loss. This is a well-known
asymmetry of handset-limited cellular links, and it is why uplink-only link budgets are
optimistic.

\paragraph{Indoor coverage.} The 12 dB urban penetration margin collapses the radius from
1580 m to 717 m (Table~\ref{tab:lb_indoor} in the calculation sheet), which is why the hospital
needs a dedicated small cell rather than relying on the outdoor macro.

\subsection{Site count and layout}
The propagation model is the log-distance law~\cite{rappaport}. A three-sector site serves one hexagonal cell of area $\frac{3\sqrt3}{2}R^2 = 2.598R^2$:

\begin{center}\small
\begin{tabular}{@{}lccc@{}}
\toprule
& $R$ & Cell area & Sites for the whole area \\
\midrule
Urban & 1580 m & 6.48 km$^2$ & $4.0/6.48 = 0.62$, i.e. the area is smaller than one cell \\
Rural & 2502 m & 16.26 km$^2$ & $25.0/16.26 = 1.54$ \\
\bottomrule
\end{tabular}
\end{center}

Geometry alone therefore suggests 1 urban and 2 rural sites. But geometry is not the only
criterion: the urban area must also carry 60 Mbit/s, and a single site must survive its failure.
A Monte-Carlo check of 30\,000 random residents against the 90\%-reliability threshold confirms
the choice (Table~\ref{tab:sweep}).

\begin{table}[H]\centering\small
\caption{Monte-Carlo coverage against site count (30\,000 random residents, 90\% reliability).}
\label{tab:sweep}
\begin{tabular}{@{}ccccc@{}}
\toprule
& \multicolumn{2}{c}{\textbf{Urban (4 km$^2$)}} & \multicolumn{2}{c}{\textbf{Rural (25 km$^2$)}} \\
\cmidrule(lr){2-3}\cmidrule(lr){4-5}
Sites & outdoor & indoor & outdoor & indoor \\
\midrule
%%SWEEP%%
\bottomrule
\end{tabular}
\end{table}

\paragraph{Reading the sweep.} One urban site already gives 99.0\% outdoor coverage, because the
4 km$^2$ area is smaller than a single 6.48 km$^2$ cell. We still build \key{two} urban sites, for
two reasons the coverage table cannot show: capacity (60 Mbit/s needs two sites at 37.5\% load
each) and survivability (a single site failure would otherwise wipe out the whole neighbourhood).
Rural needs \key{two} sites to cross 95\%; one site reaches 94.9\% and misses the target.

\begin{figure}[H]\centering
\begin{subfigure}{0.455\textwidth}\centering
\includegraphics[width=\textwidth]{figures/fig_coverage_urban.png}
\caption{Urban: 2 macros, hospital marked}
\end{subfigure}\hfill
\begin{subfigure}{0.455\textwidth}\centering
\includegraphics[width=\textwidth]{figures/fig_coverage_rural.png}
\caption{Rural: 2 macros, clinic marked}
\end{subfigure}
\caption{Layout and coverage. Colour is the predicted 900 MHz path loss; the red dashed line is
the 90\%-reliability outdoor boundary; tinted wedges are the three 120$^\circ$ sectors.}
\label{fig:layout}
\end{figure}

\section{Q3: Fading and modulation}
\label{sec:q3}

\subsection{Doppler shift at 50 km/h}
The maximum Doppler shift~\cite{rappaport} is $f_D = v f_c / c$ with $v = 50/3.6 = 13.89$ m/s:
\begin{align*}
f_{D,900}  &= \frac{13.89 \times 900\times10^{6}}{3\times10^{8}} = 41.67\ \text{Hz},\\
f_{D,1800} &= \frac{13.89 \times 1800\times10^{6}}{3\times10^{8}} = 83.33\ \text{Hz}.
\end{align*}
The coherence time $T_c\approx0.423/f_D$ is 10.2 ms at 900 MHz and 5.1 ms at 1800 MHz,
several 1 ms subframes, so the channel is \emph{slow} relative to the frame structure and
pilot-based estimation every subframe tracks it comfortably. We therefore need no fast feedback
loop, only the 2 dB margin for residual estimation error.

\subsection{RMS delay spread}
The RMS delay spread follows the standard moment formula~\cite{rappaport}. Normalising the given powers $[0,-3,-8]$ dB gives linear weights summing to
$1.0+0.501+0.158=1.659$:

\begin{center}\small
\begin{tabular}{@{}cccc@{}}
\toprule
Delay $\tau$ ($\mu$s) & Power (dB) & Linear power & Normalised weight $p_i$ \\
\midrule
0 & 0 & 1.000 & 0.6025 \\
1 & $-3$ & 0.501 & 0.3020 \\
3 & $-8$ & 0.158 & 0.0955 \\
\bottomrule
\end{tabular}
\end{center}

\begin{align*}
\bar\tau &= \textstyle\sum p_i\tau_i = 0.6025(0)+0.3020(1)+0.0955(3) = 0.5885\ \mu\text{s},\\
\overline{\tau^2} &= \textstyle\sum p_i\tau_i^2 = 0.3020(1)+0.0955(9) = 1.1614\ \mu\text{s}^2,\\
\sigma_\tau &= \sqrt{\overline{\tau^2}-\bar\tau^2}
= \sqrt{1.1614-0.3463} = \keym{0.9029\ \mu\text{s}}.
\end{align*}

The coherence bandwidth~\cite{rappaport} is $B_c\approx1/(5\sigma_\tau)=221.5$ kHz at the 50\% correlation level
and $1/(50\sigma_\tau)=22.2$ kHz at 90\%. The LTE subcarrier spacing is 15 kHz and the normal
cyclic prefix is 4.7 $\mu$s, so the RMS spread of 0.90 $\mu$s sits \emph{inside} the cyclic
prefix: ISI is absorbed by the CP and removed by the frequency-domain equaliser. It is not
negligible, though: it is 0.19 of the CP length, which is why the equaliser and a 2 dB
margin are both carried.

\begin{figure}[H]\centering
\includegraphics[width=0.34\textwidth]{figures/fig_delay.png}
\caption{The given power-delay profile. The RMS delay spread of 0.90 $\mu$s fits inside the
4.7 $\mu$s LTE cyclic prefix.}
\label{fig:delay}
\end{figure}

\subsection{BER simulation versus theory}
The simulation transmits $4\times10^{5}$ random bits per point through an AWGN channel and decides
by sign. For QPSK the symbols are Gray-mapped and $E_s/N_0=2E_b/N_0$, so the per-information-bit
error rate should coincide with BPSK, a property worth checking because a common simulation
error double-counts the symbol energy and makes QPSK look 3 dB better than it is.
Figure~\ref{fig:ber} shows agreement within Monte-Carlo scatter across the whole range.

\begin{figure}[H]\centering
\begin{subfigure}{0.475\textwidth}\centering
\includegraphics[width=\textwidth]{figures/fig_ber.png}
\caption{Simulation (markers) against theory (line).}
\end{subfigure}\hfill
\begin{subfigure}{0.37\textwidth}\centering
\includegraphics[width=\textwidth]{figures/fig_impairments.png}
\caption{Sensitivity at $E_b/N_0=10$ dB.}
\end{subfigure}
\caption{BPSK and QPSK in AWGN, and the effect of carrier-phase and symbol-timing error.}
\label{fig:ber}
\end{figure}

At $E_b/N_0=6$ dB the theory~\cite{berawgn} gives $Q(\sqrt{3.981})=Q(1.995)\approx0.0230$, i.e. a 2.3\% BER, and
the simulation lands on the same value. At 10 dB both fall to about $5\times10^{-6}$.

\paragraph{Why QPSK is chosen over BPSK.} At the same $E_b/N_0$ the two have the same BER, but
QPSK carries two bits per symbol, so it doubles spectral efficiency for free. BPSK would need
twice the bandwidth for the same throughput. The only reason to prefer BPSK is a marginally
simpler receiver, which is irrelevant at macro-cell scale.

\subsection{Receiver errors and the design}
The impairment sweep is the practically important part, because the theoretical curve assumes
perfect synchronisation.

\begin{table}[H]\centering\small
\caption{QPSK BER at $E_b/N_0=10$ dB under carrier-phase and timing error (simulated).}
\label{tab:impair}
\begin{tabular}{@{}cccc@{}}
\toprule
Phase error ($^\circ$) & BER & Timing error (\% of $T_s$) & BER \\
\midrule
%%IMPAIR_PHASE%%
\bottomrule
\end{tabular}
\end{table}

Three conclusions drive the design:

\begin{itemize}[leftmargin=1.4em,itemsep=1pt,topsep=2pt]
\item \textbf{Carrier phase is the fragile quantity.} A 30$^\circ$ residual error raises BER from
$5\times10^{-6}$ to about 2.5\%, a factor of 5000. The constellation has simply rotated into
its neighbour's decision region. So the design must budget for carrier recovery (a Costas loop or
pilot-aided tracking) and must not rely on absolute phase without pilots: this is precisely why
LTE carries cell-specific reference signals in every subframe.
\item \textbf{Timing error is more forgiving but not free.} Half a symbol of timing error
degrades BER to roughly 25\%, because the matched filter samples the zero-crossing instead of the
peak. Timing recovery (Gardner or early-late gate) is therefore mandatory, but a tenth of a
symbol costs little, which tells us the tracking loop need only be moderately tight.
\item \textbf{Design response.} Section~\ref{sec:q5} sets handover hysteresis and TTT so that
re-synchronisation happens while the signal is still strong; Section~\ref{sec:q6} chooses a codec
whose robustness lets us spend the remaining link margin on synchronisation rather than raw power.
Concretely: hold $E_b/N_0$ near 10 dB at the cell edge and keep residual phase error below
10$^\circ$.
\end{itemize}

\section{Q4: Interference control and capacity planning}
\label{sec:q4}

\subsection{Reuse factor and co-channel interference}
With $D/R=\sqrt{3N}$ and the first-tier approximation
$\frac{C}{I}=\frac{1}{i_0}\left(\frac{D}{R}\right)^{n}$, the reuse ratio is
$\sqrt{3}=1.732$ at $N=1$ and $\sqrt{9}=3.000$ at $N=3$.

\begin{table}[H]\centering\small
\caption{Estimated co-channel $C/I$ for the two reuse factors, with $i_0$ dominant interferers.}
\label{tab:reuse}
\begin{tabular}{@{}lccccc@{}}
\toprule
\textbf{Area} & $N$ & $D/R$ & \textbf{Omni} ($i_0=6$) & \textbf{3-sector} ($i_0=2$) & \textbf{6-sector} ($i_0=1$) \\
\midrule
%%REUSE%%
\bottomrule
\end{tabular}
\end{table}

\paragraph{Interpretation.} For $N=1$ and $n=3.5$ the omni $C/I$ is only 0.57 dB; the system is
interference-limited to the point of being unusable without sectorisation or fractional reuse.
Sectorising to three sectors per site raises it to 5.34 dB, because the directional antenna
rejects two-thirds of the interferers ($i_0$ falls from 6 to 2). Moving to $N=3$ almost doubles
$D/R$ and lifts the omni figure to 8.92 dB, comfortably above the 7 dB requirement. Rural, with
the gentler $n=3.0$, is worse at $N=1$ ($-0.62$ dB), a useful warning that a ``rural is
easier'' assumption is false for interference, because the slower decay means distant interferers
hurt more.

\paragraph{The trade-off in capacity terms.} $N=3$ triples the reuse distance but divides the
available spectrum by three. On one 10 MHz carrier:

\begin{center}\small
\begin{tabular}{@{}lcc@{}}
\toprule
& $N=1$ & $N=3$ \\
\midrule
Spectrum per cell & 10 MHz & 3.33 MHz \\
Spectral efficiency achievable at that $C/I$ & 2.5 bit/s/Hz & 2.5 bit/s/Hz \\
Capacity per sector & 20.0 Mbit/s & 6.67 Mbit/s \\
\bottomrule
\end{tabular}
\end{center}

So $N=1$ gives three times the capacity \emph{provided} interference is controlled another way.
This is exactly what LTE does: $N=1$ on the 1800 MHz capacity layer, with sectorisation and a
2 dB interference margin, and a logical $N=3$ split of the 900 MHz coverage layer.

\subsection{Adjacent-channel interference}
Co-channel interference is controlled by \emph{geometry} (reuse distance, antenna pattern), but
adjacent-channel interference is a \emph{filtering and linearity} problem, so it needs different
remedies and does not improve simply by moving cells further apart:

\begin{itemize}[leftmargin=1.4em,itemsep=1pt,topsep=2pt]
\item Keep a guard band between the two 10 MHz carriers and use a transmit mask with shoulders
steeper than the 3GPP TS 36.104 general-purpose mask~\cite{ts36104}; a spectrum-analyser sweep of the RRU is the
acceptance test.
\item Prevent strong local signals from overloading the receiver: a handset 2 m from a small cell
can saturate the LNA and splatter into adjacent channels. The small cells here use 1 W and a
minimum coupling-loss assumption rather than 43 dBm.
\item Aim the sector antennas' back lobes away from the neighbouring carrier's coverage: the
65$^\circ$ panels have front-to-back ratios better than 25 dB, which is what keeps ACI out of
edge users.
\end{itemize}

\subsection{Busy-hour demand and usable capacity}
Demand is given directly: $300\times0.2=60$ Mbit/s urban and $50\times0.1=5$ Mbit/s rural.

Usable capacity per sector uses a \emph{justified} spectral efficiency of 2.5 bit/s/Hz: what a
10 MHz FDD carrier with 2$\times$2 MIMO and 64-QAM, averaged over the cell (edge users on QPSK,
centre users on 64-QAM), realistically delivers; it is below the theoretical 6 bit/s/Hz for
64-QAM because of pilot and control overhead. Removing 20\% for control channels, reference
signals and guard bands:
\[
C_{\text{sector}} = \underbrace{10\ \text{MHz}\times2.5\ \text{bit/s/Hz}}_{25.0\ \text{Mbit/s}}
\times(1-0.20) = 20.0\ \text{Mbit/s}.
\]

\begin{table}[H]\centering\small
\caption{Capacity against demand. Backhaul is 1 Gbit/s per site.}
\label{tab:cap}
\begin{tabular}{@{}lccccccc@{}}
\toprule
\textbf{Area} & \textbf{Sites} & \textbf{Carriers} & \textbf{Per site} & \textbf{Total} &
\textbf{Demand} & \textbf{Load} & \textbf{Backhaul} \\
& & & Mbit/s & Mbit/s & Mbit/s & (\%) & load (\%) \\
\midrule
%%CAPACITY%%
\bottomrule
\end{tabular}
\end{table}

\paragraph{Backhaul is not the limit.} Total demand is 60 Mbit/s urban against 2 Gbit/s of
backhaul across two sites: a 3.0\% load. The binding constraints are radio capacity and coverage,
which is why the backhaul line sits inside the site price and needs no upgrade.

\subsection{Sectoring, cell splitting and small cells compared}
\begin{table}[H]\centering\small
\caption{Three ways to add capacity, compared for this network.}
\begin{tabularx}{\textwidth}{@{}p{2.4cm}XXX@{}}
\toprule
& \textbf{Sectoring} & \textbf{Cell splitting} & \textbf{Small cells} \\
\midrule
Mechanism & three 65$^\circ$ panels per site instead of omni & more sites at smaller $R$ & low-power nodes under macro coverage \\
Capacity gain & about 3 times per site; $C/I$ rises 4.8 dB urban & capacity per area scales as $1/R^2$ & large locally, only where deployed \\
Cost & low: antennas on an existing mast & high: a whole BDT 6 M site each & lowest: BDT 1 M per node \\
Interference cost & lower ($i_0$: 6 to 2) & higher: more reuse cells closer together & needs careful indoor/outdoor coordination \\
Handovers & unchanged & increases, each a drop opportunity & many, between layers \\
Best for & both areas, already adopted & growth beyond the budget & the hospital and clinic exactly \\
\bottomrule
\end{tabularx}
\end{table}

\paragraph{Selection.} \key{Urban: sectoring plus the 1800 MHz $N=1$ capacity layer, with small
cells reserved for the hospital.} Sectoring is nearly free and already gives the 4.8 dB of $C/I$
needed to make $N=1$ viable; the 1800 MHz carrier triples per-site capacity to 80 Mbit/s.
\key{Rural: sectoring plus the 900 MHz coverage layer, no small cells outside the clinic.}
Splitting rural cells would waste money: two sites already run at 12.5\% load, so more capacity is
not the problem; reach is, and that is bought with tower height and the 900 MHz band, not with
more sites.

\section{Q5: Handover improvement}
\label{sec:q5}

\subsection{The model}
A handset drives at 50 km/h~\cite{handover} (13.89 m/s) from inside site A, across the boundary midway between
the two sites, into site B. Site separation is $2R=3.16$ km, so the midpoint of the route is
exactly the boundary. Every 10 ms the simulator evaluates
\[
P_r(d) = \underbrace{P_t+G_t+G_r-L_{\text{body}}}_{EIRP=35.5\ \text{dBm}}
-\left[PL(d_0)+10n\log_{10}\frac{d}{d_0}\right] + X_\sigma + F,
\]
with correlated log-normal shadowing $X_\sigma$ ($\sigma=8$ dB, 50 m decorrelation) and a
Jakes-sum Rayleigh term $F$ at $f_D=41.67$ Hz. A 200 ms measurement window averages the fast
fading, matching how a handset actually reports. The handover rule is the one from the course
notes:
\[
P_{\text{target}} > P_{\text{serving}} + H \quad\text{sustained for the time-to-trigger (TTT)}.
\]

\begin{figure}[H]\centering
\includegraphics[width=0.60\textwidth]{figures/fig_handoff.png}
\caption{Received-signal traces for the two settings. The black line is the serving-cell level.
(a) $H=3$ dB, no TTT: the traces cross repeatedly near the boundary and the call is handed back
and forth. (b) $H=6$ dB, TTT $=1$ s: one decisive handover.}
\label{fig:handoff}
\end{figure}

\subsection{Comparing the two settings}
\begin{table}[H]\centering\small
\caption{Handover performance over 30 random channel realisations per setting.}
\label{tab:handoff}
\begin{tabular}{@{}ccccccc@{}}
\toprule
$H$ (dB) & TTT (s) & Handovers & 95th pct. & Ping-pongs & Degraded (\%) & RSRP at HO (dBm) \\
\midrule
%%HANDOFF_STATS%%
\bottomrule
\end{tabular}
\end{table}

\paragraph{The trade-off is monotone in the table.} At $H=2$ dB with no TTT the handset performs about
32 handovers per crossing with 23 ping-pongs, so the network is thrashing. Raising to $H=9$ dB with
a 2 s TTT cuts that to 1.7 handovers and near-zero ping-pongs, but degraded time climbs to
5.69\%. That is the classical conflict:

\begin{itemize}[leftmargin=1.4em,itemsep=1pt,topsep=2pt]
\item \textbf{Too little hysteresis} gives \emph{repeated switching}: signalling load,
interruption on every hard handover, and a higher drop probability simply because there are more
chances to fail.
\item \textbf{Too much hysteresis} gives \emph{late handover}: the call stays on a cell that has
already faded, so the terminal transmits at maximum power into a poor channel and can hit
radio-link failure before the handover completes.
\end{itemize}

The optimum is where the ping-pong count has already collapsed but the degraded-time fraction has
not yet risen. Table~\ref{tab:handoff} puts that region at $H=3$-$4$ dB with TTT around 1 s.

\subsection{Recommended settings}
\begin{itemize}[leftmargin=1.4em,itemsep=1pt,topsep=2pt]
\item \textbf{Urban: $H=3$ dB, TTT $=0.5$ s.} Cells are small (1.58 km) and users cross many
boundaries, so handovers must be prompt. 0.5 s of TTT is long enough to reject Rayleigh dips (the
fade decorrelates in about 10 ms, so a 0.5 s window spans many independent fades) but short enough
to hand over before the $-85$ dBm threshold. Ping-pongs fall from about 23 to 3.
\item \textbf{Rural: $H=6$ dB, TTT $=1$ s.} Cells are large (2.5 km), so crossings are rare and
there is time to be conservative. A larger hysteresis suppresses the ping-pong that would
otherwise occur across flat rural terrain, where the two site levels stay within a few dB for
hundreds of metres, and 1 s of TTT costs almost nothing because the handset spends only 0.02\% of
the crossing below threshold.
\end{itemize}

\subsection{How dropped calls would be measured}
Define a dropped call operationally, then instrument it:
\begin{enumerate}[leftmargin=1.4em,itemsep=1pt,topsep=2pt]
\item \textbf{RRC connection abnormal release} with a cause other than normal release and no
subsequent re-establishment within a short window. This is the primary KPI, from MME/eNB counters. Call
drop rate equals abnormal releases divided by successfully established connections, per cell and
per area, reported hourly and at busy hour.
\item \textbf{Handover-specific drops:} handover failures (target preparation failure,
re-establishment failure) divided by handover attempts. This isolates the late-handover mechanism
from coverage drops.
\item \textbf{Radio-link failure (RLF) count} per cell, together with the RSRP and SINR at the
moment of failure. If failures cluster just before a boundary, hysteresis is too large; if they
cluster mid-cell, it is coverage.
\item \textbf{End-to-end corroboration:} a synthetic call generator plus a drive test on the two
crossing routes, logging RSRP, SINR, handover events and gaps. This catches failures that network
counters attribute elsewhere.
\end{enumerate}

\section{Q6: Speech coding and voice capacity}
\label{sec:q6}

\subsection{PCM bit rate}
For 8-bit PCM sampled at 8 kHz~\cite{g711}:
\[
R_b = f_s\times n = 8000\times8 = \key{64000\ \text{bit/s} = 64\ \text{kbit/s}}.
\]
The quantisation signal-to-noise ratio for a full-scale sinusoid is
$6.02n+1.76 = 6.02(8)+1.76 = 49.92$ dB, and the ideal minimum baseband bandwidth is
$R_b/2 = 32$ kHz.

\subsection{Comparison of coding methods}
\begin{table}[H]\centering\small
\caption{Speech coding techniques compared. The star marks the recommendation.}
\label{tab:codecs}
\begin{tabularx}{\textwidth}{@{}p{4.3cm}cp{1.9cm}p{2.0cm}p{1.7cm}X@{}}
\toprule
\textbf{Method} & \textbf{Bit rate} & \textbf{Bandwidth} & \textbf{Quality} & \textbf{Complexity} & \textbf{Typical use} \\
& kbit/s & & & & \\
\midrule
%%CODECS%%
\bottomrule
\end{tabularx}
\end{table}

\paragraph{Why the ordering is what it is.} PCM is a \emph{waveform} coder: it transmits the sample
itself, so quality is high and complexity almost zero, but 64 kbit/s is expensive over radio.
DPCM transmits the \emph{prediction error}; because adjacent speech samples correlate strongly,
4 bits suffice where PCM needs 8, halving the rate at a small quality cost. Delta modulation goes
further, to 1 bit per sample, by transmitting only the \emph{direction} of change, hence low
complexity but slope-overload and granular noise. LPC and the vocoders model the \emph{vocal
tract} instead of the waveform, so 2.4-13 kbit/s is achievable, at the cost of significant DSP
and a more synthetic timbre. The AMR-NB family used in GSM and VoLTE~\cite{amr,ts26071} sits at the sweet spot:
12.2 kbit/s, near-toll quality, and adaptive between modes.

\paragraph{Recommendation.} \key{AMR-NB at 12.2 kbit/s (ACELP, an LPC-class parametric coder)},
with EVS as the upgrade path where handsets support it. It is 5.2 times more efficient than PCM,
delivers 3.9 MOS against PCM's 4.3 at 14\% of the rate, and is the standard VoLTE codec, so it
interoperates with the existing core. Delta modulation is rejected because its 2.8 MOS quality is
audible on a mobile link, and LPC-10 is rejected because 2.6 MOS synthetic speech is unacceptable
for a service the regulator will monitor.

\subsection{Offered voice traffic and Erlang-B dimensioning}
With 0.02 Erlang per subscriber (a busy-hour average busy fraction), offered traffic is
\[
A_{\text{urban}} = 0.02\times10\,000 = 200\ \text{E},\qquad
A_{\text{rural}} = 0.02\times2\,000 = 40\ \text{E}.
\]
Erlang B~\cite{erlang}, $B(A,N)=\frac{A^N/N!}{\sum_{k=0}^{N}A^k/k!}$, is solved for the smallest $N$ with
$B\leq2\%$.

\begin{table}[H]\centering\small
\caption{Erlang-B dimensioning for voice at 2\% blocking.}
\label{tab:voice}
\begin{tabular}{@{}lccccccc@{}}
\toprule
& Subscribers & Offered $A$ (E) & Channels $N$ & $B$ at $N$ & $B$ at $N{-}1$ & Utilisation & Per sector \\
\midrule
%%VOICE%%
\bottomrule
\end{tabular}
\end{table}

\paragraph{Reading the table.} Urban needs 214 channels; at 213 the blocking rises to 2.18\% and
misses the target, so 214 is the true minimum, a good sanity check that the dimensioning is not
over-provisioned. Rural needs 50 channels. Spread over two three-sector sites this is 35.7 and 8.3
channels per sector respectively, both trivially available inside a 20 Mbit/s VoLTE sector.

\paragraph{Why packet voice needs extra overhead accounting.} Circuit-switched voice consumes
exactly one 64 kbit/s channel. Packet voice does not: the speech payload travels inside
RTP/UDP/IP headers that are pure overhead. With 12.2 kbit/s AMR in 20 ms frames the payload is
30.5 bytes per frame, while RTP(12)+UDP(8)+IPv4(20) and the LTE PDCP/RLC/MAC headers add 46
bytes:

\begin{center}\small
\begin{tabular}{@{}lccc@{}}
\toprule
& Payload (B) & Headers (B) & Wire rate (kbit/s) \\
\midrule
No header compression & 30.5 & 46 & 30.6 \\
With ROHC (3 B) & 30.5 & 9 & 15.8 \\
\bottomrule
\end{tabular}
\end{center}

Without compression the 12.2 kbit/s codec actually occupies 30.6 kbit/s on the air, a 151\%
overhead, so more than half the radio capacity is headers. Robust Header Compression (ROHC) cuts
this to 29.5\% and 15.8 kbit/s. This is why VoLTE capacity planning must be done on the \emph{wire}
rate, not the codec rate, and why semi-persistent scheduling matters: it removes the per-packet
grant signalling that would otherwise dominate for small voice packets.

\section{Q7: Defending the final design}
\label{sec:q7}

\subsection{The costed plan}
Every line below is a real equipment class; the full bill of materials with vendor and standards
links is in Appendix~\ref{app:bom}. Unit costs are the assignment's teaching values; the linked
datasheets justify the \emph{technical} choice, not the price.

\begin{table}[H]\centering\small
\caption{Budget roll-up against the BDT 30.0 M cap (full BOM with links: Appendix~\ref{app:bom}).}
\label{tab:budget}
\begin{tabular}{@{}lcccc@{}}
\toprule
\textbf{Line} & \textbf{Item} & \textbf{Qty} & \textbf{Unit (BDT M)} & \textbf{Total (BDT M)} \\
\midrule
%%BUDGET%%
\midrule
\multicolumn{4}{@{}l}{\textbf{Capital cost}} & \textbf{28.8} \\
\multicolumn{4}{@{}l}{10\% contingency} & 2.88 \\
\multicolumn{4}{@{}l}{\textbf{Total including contingency}} & \textbf{31.68} \\
\multicolumn{4}{@{}l}{Budget cap} & 30.00 \\
\multicolumn{4}{@{}l}{\textbf{Overrun with contingency}} & \textbf{+1.68} \\
\bottomrule
\end{tabular}
\end{table}

\paragraph{Backhaul capacity.} Each site is provisioned with 1 Gbit/s (licensed 18/23 GHz
microwave for the rural hops, GPON fibre where duct exists). Against 60 Mbit/s of urban
busy-hour demand that is a 3.0\% load; backhaul is therefore explicitly \emph{not} the binding
constraint, and no backhaul upgrade is costed.

\subsection{Two alternatives compared}
\begin{table}[H]\centering\small
\caption{Alternatives against coverage, capacity, cost and emergency continuity.}
\begin{tabularx}{\textwidth}{@{}p{2.9cm}XXX@{}}
\toprule
& \textbf{A: recommended} & \textbf{B: coverage-first} & \textbf{C: capacity-first} \\
\midrule
Urban & 2 macros, 900+1800 MHz, hospital small cell & 1 macro + 4 small cells & 2 macros + 3 small cells \\
Rural & 2 macros 900 MHz + clinic small cell & 2 macros 900 MHz, no small cell & 3 macros, 1800 MHz added \\
Outdoor coverage & U 99.0\% / R 97.8\% & U 99.0\% / R 97.8\% & U 99.9\% / R 99.0\% \\
Busy-hour load & U 37.5\% / R 12.5\% & U 75\% (single site) & U 25\% / R 8\% \\
Site-failure tolerance & survives (75\% on survivor) & \textbf{fails}: one site is the area & survives easily \\
Hospital / clinic & both on dedicated small cells & hospital only & both \\
Capital cost & \textbf{BDT 28.8 M} & about BDT 24.8 M & BDT 35.6 M \\
Within budget? & yes (4.0\% headroom) & yes & \textbf{no, +5.6 M} \\
\bottomrule
\end{tabularx}
\end{table}

\paragraph{Why B is rejected despite being cheapest.} A single urban macro looks attractive:
coverage is identical and it saves a full BDT 6 M site. But it runs at 75\% load with \emph{no
redundancy at all}: any single failure removes the whole neighbourhood, and the hospital small
cell's donor goes with it. It also cannot absorb the 50\% growth test. Saving BDT 4 M by making an
entire area a single point of failure is not a defensible engineering trade.

\paragraph{Why C is rejected.} It is technically the strongest, with more capacity and the best
coverage, but it costs \textbf{BDT 35.6 M, i.e. BDT 5.6 M (18.7\%) over the cap}. It is recorded
here as the target state for the staged upgrade rather than as a candidate for this budget.

\subsection{Test 1: one site failure}
Remove one of the two urban macro sites; the survivor must carry the entire 60 Mbit/s demand from
its own 80 Mbit/s of capacity:
\[
\text{load} = \frac{60}{80} = 75.0\% \;(<100\%) \;\Rightarrow\; \text{service held, degraded.}
\]
The area degrades (users near the lost site fall back to a more distant cell and see lower
throughput) but no user is denied service, and the hospital small cell keeps working because its
coverage is now provided by the survivor. Under the \emph{combined} worst case of one site failure
\emph{and} 50\% growth, the survivor must carry 90 Mbit/s against 80 Mbit/s, i.e. 112.5\%, a
10 Mbit/s shortfall. That is quantified honestly rather than hidden: it is the trigger condition
for the Phase-2 site.

\subsection{Test 2: 50\% traffic increase}
\begin{center}\small
\begin{tabular}{@{}lccc@{}}
\toprule
& Demand (Mbit/s) & Capacity (Mbit/s) & Load (\%) \\
\midrule
Urban, +50\% & 90.0 & 160.0 & 56.2 \\
Rural, +50\% & 7.5 & 40.0 & 18.7 \\
\bottomrule
\end{tabular}
\end{center}

Both areas absorb 50\% growth without new hardware (56.2\% and 18.7\% load). The rural area has so
much headroom that it could absorb 400\% growth before needing a third carrier, which is the
quantitative argument for not building the rural 1800 MHz layer now.

\subsection{Resolving stakeholder priorities}
\begin{table}[H]\centering\small
\caption{Conflicting stakeholder priorities and the resolution adopted.}
\begin{tabularx}{\textwidth}{@{}p{2.5cm}p{3.5cm}X@{}}
\toprule
\textbf{Stakeholder} & \textbf{Priority} & \textbf{How it is resolved in the design} \\
\midrule
Operator & lowest cost per covered user; capex discipline & Prioritise sectoring over new sites (nearly free capacity), one shared DC plant per site rather than per radio, and defer the third urban site to Phase 2. \\
Community & reliable everyday service; visible fairness & Both areas get two macros so neither is obviously favoured; rural gets the more expensive 24-hour energy solution because flood risk is higher there. \\
Hospital / clinic & continuity during an outage, above all & Dedicated small cells on dedicated DC plants; the 8-hour urban outage does not touch the clinic's 24-hour one. Continuity is \emph{hard-allocated}, not hoped for. \\
Regulator & blocking and coverage targets met & Dimensioned to 1.99\% blocking and 97.8\% worst-area coverage; both are documented so they can be audited. \\
\bottomrule
\end{tabularx}
\end{table}

The principle: \emph{protect the irreducible requirement, trade on the elastic one}. Hospital
continuity is irreducible and hard-allocated; capacity is elastic and met with spectrum and
sectorisation before steel; cost is traded last.

\subsection{The budget shortfall and a staged upgrade}
The base plan is BDT 28.8 M, which fits with BDT 1.2 M of headroom. But a prudent 10\%
contingency is BDT 2.88 M, taking the total to BDT 31.68 M, which is \key{BDT 1.68 M over the cap}.
That is stated as genuine infeasibility, not absorbed by wishful thinking.

\begin{table}[H]\centering\small
\caption{Staged upgrade plan with quantified shortfall.}
\label{tab:stages}
\begin{tabularx}{\textwidth}{@{}p{0.75cm}p{4.7cm}Xp{1.85cm}@{}}
\toprule
\textbf{Stage} & \textbf{Scope} & \textbf{Justification and trigger} & \textbf{Cost} \\
\midrule
1 & 2 urban macros, 2 rural macros, 2 small cells, 8 h urban backup, 24 h rural backup & Now. Meets every stated target. & BDT 28.8 M \\
2 & Fund the contingency from first-year opex & Now. The 1.68 M overrun is closed by funding the contingency reserve from operating cash flow, normal practice for a plan with 1.2 M of unallocated headroom. & +0 (timing) \\
3 & Third urban macro site and 8 h backup & When busy-hour load exceeds 70\%, or a site fails during growth. Restores N+1 and covers the 10 Mbit/s shortfall quantified above. & BDT 6.4 M \\
4 & Rural 1800 MHz carrier & When rural demand exceeds 20 Mbit/s. Not needed until rural demand grows fourfold; adding it early would waste BDT 1 M of spectrum capex. & BDT 1.0 M \\
\bottomrule
\end{tabularx}
\end{table}

So the honest answer to ``does it fit?'' is: \textbf{the engineering fits at BDT 28.8 M, but a
fully risk-reserved plan does not fit by BDT 1.68 M}, and the third site would not fit by BDT
5.2 M. We deliver Stages 1-2 inside the cap and hold Stage 3 as a measured, triggered commitment.

\subsection{Target-compliance table}
All eight stated targets are met; the full table with the checking method for each is
Table~\ref{tab:compliance} in Appendix~\ref{app:calc}.

\section{Limitations and what would change the answer}
\begin{itemize}[leftmargin=1.4em,itemsep=1pt,topsep=2pt]
\item The log-distance model with a single exponent ignores terrain. Kurigram is flood-prone and
flat; a real design would use ITU-R P.1546~\cite{p1546} with a digital elevation model plus clutter data, which
could move the rural radius by $\pm15\%$.
\item Sectors are modelled as ideal 65$^\circ$ patterns. Real patterns have back lobes that raise
$i_0$ above 2 and cut the $N=1$ $C/I$ by 1-2 dB.
\item The Erlang-B model assumes Poisson arrivals and no handover priority. Real networks hold
channels for handover, which raises blocking slightly above the computed 1.99\%.
\item Prices are the assignment's teaching unit costs. A real procurement would add civil works,
spectrum fees, site rental and opex, none of which are inside the BDT 30 M cap.
\item The coverage prediction is a model, not a measurement. The compliance table therefore
specifies a drive test as the acceptance criterion, so a model error becomes a field finding
rather than a silent failure.
\end{itemize}

\renewcommand{\refname}{References}
\begin{multicols}{2}
\begin{thebibliography}{99}
\setlength{\itemsep}{0pt}
\footnotesize

\bibitem{ts36104} 3GPP, ``E-UTRA: Base Station (BS) radio transmission and reception,'' TS 36.104.
[Online]. Available: \url{https://www.3gpp.org/ftp/Specs/archive/36_series/36.104/}

\bibitem{ts36101} 3GPP, ``E-UTRA: User Equipment (UE) radio transmission and reception,''
TS 36.101. [Online]. Available:
\url{https://www.3gpp.org/ftp/Specs/archive/36_series/36.101/}

\bibitem{ts36942} 3GPP, ``E-UTRA: Radio Frequency (RF) system scenarios,'' TS 36.942. [Online].
Available: \url{https://www.3gpp.org/ftp/Specs/archive/36_series/36.942/}

\bibitem{ts26071} 3GPP, ``Mandatory speech codec speech processing functions; AMR speech codec;
General description,'' TS 26.071. [Online]. Available:
\url{https://www.3gpp.org/ftp/Specs/archive/26_series/26.071/}

\bibitem{ts23228} 3GPP, ``IP Multimedia Subsystem (IMS); Stage 2,'' TS 23.228. [Online].
Available: \url{https://www.3gpp.org/ftp/Specs/archive/23_series/23.228/}

\bibitem{p1546} ITU, ``Method for point-to-area predictions for terrestrial services in the
frequency range 30 MHz to 3000 MHz,'' Recommendation ITU-R P.1546. [Online]. Available:
\url{https://www.itu.int/rec/R-REC-P.1546/en}

\bibitem{p2109} ITU, ``Prediction of building entry loss,'' Recommendation ITU-R P.2109.
[Online]. Available: \url{https://www.itu.int/rec/R-REC-P.2109/en}

\bibitem{p837} ITU, ``Characteristics of precipitation for propagation modelling,''
Recommendation ITU-R P.837. [Online]. Available:
\url{https://www.itu.int/rec/R-REC-P.837/en}

\bibitem{p526} ITU, ``Propagation by diffraction,'' Recommendation ITU-R P.526. [Online].
Available: \url{https://www.itu.int/rec/R-REC-P.526/en}

\bibitem{g711} ITU, ``Pulse code modulation (PCM) of voice frequencies,'' Recommendation
ITU-T G.711. [Online]. Available: \url{https://www.itu.int/rec/T-REC-G.711/en}

\bibitem{g722} ITU, ``7 kHz audio-coding within 64 kbit/s,'' Recommendation ITU-T G.722.
[Online]. Available: \url{https://www.itu.int/rec/T-REC-G.722/en}

\bibitem{g168} ITU, ``Digital network echo cancellers,'' Recommendation ITU-T G.168. [Online].
Available: \url{https://www.itu.int/rec/T-REC-G.168/en}

\bibitem{rappaport} T. S. Rappaport, \emph{Wireless Communications: Principles and Practice},
2nd ed. Upper Saddle River, NJ, USA: Prentice Hall, 2002.

\bibitem{lee} W. C. Y. Lee, \emph{Mobile Cellular Telecommunications: Analog and Digital
Systems}, 2nd ed. New York, NY, USA: McGraw-Hill, 1995.

\bibitem{erlang} Westbay Engineers, ``Erlang B calculator.'' [Online]. Available:
\url{https://www.erlang.com/calculator/erlb/}

\bibitem{berawgn} MathWorks, ``BERAWGN: theoretical bit error rate in AWGN.'' [Online].
Available: \url{https://www.mathworks.com/help/comm/ref/berawgn.html}

\bibitem{itustats} ITU, ``ITU-D statistics.'' [Online]. Available:
\url{https://www.itu.int/en/ITU-D/Statistics/Pages/stat/default.aspx}

\bibitem{worldbank} The World Bank, ``Individuals using the Internet (\% of population),
Bangladesh.'' [Online]. Available:
\url{https://data.worldbank.org/indicator/IT.NET.USER.ZS?locations=BD}

\bibitem{gsma} GSMA, ``GSMA Intelligence: mobile economy data.'' [Online]. Available:
\url{https://www.gsmaintelligence.com/}

\bibitem{opensignal} Opensignal, ``Mobile experience reports.'' [Online]. Available:
\url{https://www.opensignal.com/}

\bibitem{ookla} Ookla, ``Speedtest Global Index.'' [Online]. Available:
\url{https://www.ookla.com/}

\bibitem{handover} A. R. Mishra, \emph{Advanced Cellular Network Planning and Optimisation:
2G/2.5G/3G\ldots Evolution to 4G}. Chichester, U.K.: Wiley, 2007.

\bibitem{amr} 3GPP, ``AMR speech codec; General description,'' TS 26.071. [Online]. Available:
\url{https://www.3gpp.org/ftp/Specs/archive/26_series/26.071/}

\end{thebibliography}
\end{multicols}

\appendix
\section{Calculation sheet}
\label{app:calc}

\subsection{Assumptions and justification}
\begin{table}[H]\centering\small
\caption{Every design value the assignment did not supply, with justification.}
\label{tab:assume}
\begin{tabularx}{\textwidth}{@{}p{3.5cm}p{1.9cm}X@{}}
\toprule
\textbf{Parameter} & \textbf{Value used} & \textbf{Justification} \\
\midrule
Reference distance $d_0$ & 100 m & Given. Free space is valid at $d_0$ for a mast-height antenna. \\
Reference loss $PL(d_0)$, 900 MHz & 71.5 dB & Free space: $20\log_{10}(4\pi d_0 f/c)$ with $d_0=100$ m, $f=900$ MHz. \\
Reference loss $PL(d_0)$, 1800 MHz & 77.6 dB & Same formula at 1800 MHz: the expected 6.02 dB offset. \\
Path-loss exponent $n$ & 3.5 urban / 3.0 rural & Given, and inside the ranges quoted in the course notes. \\
BS antenna gain & 16 dBi at 900, 18 dBi at 1800 & 65$^\circ$ three-sector panel; the 2 dB difference follows the shorter wavelength. \\
Handset gain & 0 dBi & Omnidirectional handset; body loss carried separately. \\
BS noise figure & 5 dB & Typical macro RRU. \\
Handset noise figure & 7 dB & Typical mass-market handset front end. \\
Feeder loss & 0.5 dB & RRU mounted at the mast head, so only a short jumper is lossy. \\
Body / head loss & 3 dB & Standard handset body-loss allowance. \\
Receiver $C/N$ requirement & 7 dB & Derived in Section~\ref{sec:q2}, not assumed. \\
Shadow margin (urban / rural) & 10 dB / 8 dB & 90\%-of-locations edge reliability with $\sigma=8$ dB urban, 6 dB rural. \\
Indoor margin (urban / rural) & 12 dB / 8 dB & Multi-wall concrete urban against single-wall brick rural, consistent with ITU-R P.2109 classes~\cite{p2109}. \\
Interference margin & 2 dB & Allows for a fractional-reuse or $N=1$ capacity layer. \\
Clutter margin (rural) & 2 dB & Foliage and scattered village structures. \\
Spectral efficiency & 2.5 bit/s/Hz & 10 MHz FDD, 2$\times$2 MIMO, 64-QAM averaged over the cell. \\
Frame overhead & 20\% & Control channels, reference signals and guard bands. \\
Backhaul per site & 1 Gbit/s & Licensed 18/23 GHz microwave or GPON fibre. \\
Busy-hour load target & $\leq$70\% & Headroom for growth and for one site failure. \\
\bottomrule
\end{tabularx}
\end{table}

\subsection{Worked link budget: urban 900 MHz uplink, step by step}
\begin{center}\small
\begin{tabular}{@{}lrl@{}}
\toprule
\textbf{Step} & \textbf{Value} & \textbf{Working} \\
\midrule
Handset EIRP & 20.0 dBm & $P_t + G_t - L_{\text{body}} = 23 + 0 - 3$ \\
Path loss allowed at BS & 127.5 dB & $\text{EIRP} + G_{bs} - L_{\text{feeder}} - S_{bs} = 20 + 16 - 0.5 + 92.0$ \\
Thermal noise, 10 MHz & $-99.0$ dBm & $-174 + 70 + 5$ \\
BS sensitivity & $-92.0$ dBm & $N_{bs} + C/N = -99.0 + 7.0$ \\
Total margin & 12.0 dB & shadow 10 + interference 2 \\
Net budget & 115.5 dB & $127.5 - 12.0$ \\
Log-distance inversion & 1801.7 m & $d_0\,10^{(115.5-71.5)/35}$ \\
Downlink radius & 1579.5 m & $d_0\,10^{(113.5-71.5)/35}$ \\
\textbf{Design radius} & \textbf{1579.5 m} & \textbf{downlink limits} \\
\bottomrule
\end{tabular}
\end{center}

\subsection{Site-count arithmetic}
\[
A_{\text{cell}} = \frac{3\sqrt3}{2}R^2 = 2.598R^2,\qquad
N_{\text{sites}} = \frac{A_{\text{area}}}{2.598R^2}.
\]
Urban: $2.598(1.5795)^2 = 6.48$ km$^2$, so $4.0/6.48 = 0.62$ sites by geometry (2 built for
capacity and redundancy). Rural: $2.598(2.502)^2 = 16.26$ km$^2$, so $25.0/16.26 = 1.54$, rounded
up to 2.

\subsection{Indoor link budgets}
%%LINK_INDOOR%%

\subsection{Delay-spread moments}
\[
\bar\tau = \sum_i p_i\tau_i,\qquad
\overline{\tau^2} = \sum_i p_i\tau_i^2,\qquad
\sigma_\tau = \sqrt{\overline{\tau^2}-\bar\tau^2}
\]
with $p_i = P_i/\sum_j P_j$ from the linear powers $[1.0,\ 0.501,\ 0.158]$, giving
$\bar\tau=0.5885$ $\mu$s, $\overline{\tau^2}=1.1614$ $\mu$s$^2$ and
$\sigma_\tau=0.9029$ $\mu$s.

\subsection{Erlang-B arithmetic}
\[
B(A,N) = \frac{A^N/N!}{\sum_{k=0}^{N} A^k/k!},\qquad
B(200,213) = 2.18\%,\quad B(200,214) = 1.99\%.
\]
The smallest $N$ meeting 2\% is therefore 214 channels for the urban area, and
$B(40,49)=2.38\%$, $B(40,50)=1.87\%$ gives 50 channels for the rural area.

\subsection{VoLTE overhead arithmetic}
Speech payload per 20 ms frame: $12\,200\times0.020/8 = 30.5$ bytes. RTP+UDP+IPv4 $=40$ bytes
plus PDCP+RLC+MAC $=6$ bytes gives 46 bytes of header. Wire rate without compression
$=(30.5+46)\times8/0.020=30.6$ kbit/s; with ROHC (3 bytes) it is $15.8$ kbit/s, i.e. 77\% payload
efficiency.

\subsection{Target-compliance evidence}
\begin{table}[H]\centering\footnotesize
\caption{Every stated target with the evidence that checks it.}
\label{tab:compliance}
\begin{tabularx}{\textwidth}{@{}p{0.6cm}p{4.0cm}p{2.2cm}p{3.0cm}p{0.9cm}X@{}}
\toprule
& \textbf{Target} & \textbf{Required} & \textbf{Achieved} & \textbf{Status} & \textbf{How it is checked} \\
\midrule
%%TARGETS%%
\bottomrule
\end{tabularx}
\end{table}

\section{Hardware bill of materials with vendor links}
\label{app:bom}
Every line is a real equipment class. Links are to vendor product pages or the governing 3GPP / ITU
specification, and each was checked to resolve. Unit costs are the assignment's teaching values.

\begin{table}[H]\centering\footnotesize
\caption{Hardware bill of materials: specification, rationale and references.}
\begin{tabularx}{\textwidth}{@{}p{0.65cm}p{5.8cm}cp{1.35cm}p{4.3cm}@{}}
\toprule
\textbf{ID} & \textbf{Item and specification} & \textbf{Qty} & \textbf{Unit cost} & \textbf{References} \\
\midrule
%%HARDWARE%%
\bottomrule
\end{tabularx}
\end{table}

\subsection{Supporting software and standards}

\begin{table}[H]\centering\footnotesize
\caption{Tools and standards used in the analysis.}
\begin{tabularx}{\textwidth}{@{}p{4.3cm}Xp{6.2cm}@{}}
\toprule
\textbf{Item} & \textbf{Why it is used} & \textbf{References} \\
\midrule
%%SOFTWARE%%
\bottomrule
\end{tabularx}
\end{table}

\section{Complete link register}
Every hardware and standards link used in this report, collected for verification.

{\scriptsize
\begin{longtable}{@{}p{0.65cm}p{4.4cm}p{9.7cm}@{}}
\toprule
\textbf{Ref} & \textbf{Resource} & \textbf{URL} \\
\midrule
\endhead
%%ALLLINKS%%
\bottomrule
\end{longtable}
}

\clearpage
\section{Cisco Packet Tracer network model}
\label{app:pt}

The logical design in Section~\ref{sec:q7} was also built as a runnable Cisco Packet Tracer
topology, so that the access, aggregation and core layers can be inspected as an end-to-end
packet path rather than as a costing table alone. Figure~\ref{fig:pt} is a capture of that
model. The topology file itself (\texttt{Cellular\_PT.pkt}) is supplied alongside this report
and opens in Packet Tracer 8.x; Table~\ref{tab:ptinv} lists every device in it.

\begin{figure}[H]\centering
\includegraphics[width=\textwidth]{figures/fig_pt_topology.png}
\caption{Packet Tracer topology: two access rows per area (urban Dhaka, rural Kurigram),
aggregation, and a shared EPC core with HSS, SGW, PGW, MME, IMS, plus firewall and ISP egress.}
\label{fig:pt}
\end{figure}

\subsection{How the model maps onto the design}
The simulator stands in for the radio layer rather than reproducing it. Packet Tracer has no
LTE physical layer, so the four \texttt{Cell-Tower} objects represent the four macro sites of
Table~\ref{tab:budget} (two urban, two rural), and each eNB is backhauled through a
\texttt{2960} access switch and a \texttt{Central-Office-Server} to the aggregation router.
The core follows the same EPC decomposition used in the Q4 and Q7 reasoning: MME for control
plane, SGW/PGW for the user plane, HSS for subscriber data, IMS for VoLTE, and an edge firewall
in front of ISP egress.

\begin{table}[H]\centering\footnotesize
\caption{Device inventory of the Packet Tracer model.}
\label{tab:ptinv}
\begin{tabularx}{\textwidth}{@{}p{5.3cm}p{0.8cm}X@{}}
\toprule
\textbf{Device} & \textbf{Qty} & \textbf{Role in the model} \\
\midrule
\texttt{Cell-Tower} CT-DHK-1/2, CT-KUR-1/2 & 4 & Macro sites, matching the four BDT 6.0 M sites \\
\texttt{2911} eNB-DHK-1/2, eNB-KUR-1/2 & 4 & Baseband/EPC-facing node per site \\
\texttt{2960} SW-DHK-1/2, SW-KUR-1/2 & 4 & Site access switching \\
Central-Office-Server CO-DHK/KUR & 4 & Local service/central-office host \\
\texttt{2911} AGG-DHK-1/2, AGG-KUR & 3 & Aggregation toward the core \\
\texttt{3650} CORE-SW & 1 & Core switch fabric \\
\texttt{2911} CORE-MME/SGW/PGW/IMS & 4 & EPC and IMS functions \\
SRV-HSS/DHCP/NMS/INTRA/WWW & 5 & Subscriber, addressing, management, services \\
\texttt{5505} FW-DHK, \texttt{2911} ISP-RTR & 2 & Security and external egress \\
\texttt{7960} IP-HOSP-01/02, IP-CLINIC-01 & 3 & Hospital and clinic voice terminals \\
AP-HOSP, AP-CLINIC & 2 & Indoor Wi-Fi, standing in for small cells \\
UE-DHK-01..03, UE-KUR-01..04 & 7 & Mobile subscriber endpoints \\
TAB-HOSP-01/02, TAB-CLINIC-01 & 3 & Indoor data clients \\
SIM-ATTACKER & 1 & Traffic generator used to drive load \\
\bottomrule
\end{tabularx}
\end{table}

\subsection{What the model does and does not show}
The capture reproduces the structural claim of Q7, that both areas share one core and one
budget, and it makes the hospital and clinic attachment points explicit. Two qualifications are
recorded. First, the red link indicators on \texttt{IP-HOSP-01}, \texttt{IP-HOSP-02} and
\texttt{IP-CLINIC-01} in Figure~\ref{fig:pt} mark access links that are down in the saved state,
so continuity of the protected terminals is topology intent, not a passing test. Second, the
BOM small cells appear as Wi-Fi access points, since Packet Tracer models no LTE small cell; no
BER, SINR, reuse or Erlang result here is derived from the simulator, and the numerical results
remain those of Appendices~\ref{app:calc} and~\ref{app:bom}.

\end{document}